(1) 由Faraday与Ampère–Maxwell方程，\(\nabla\times\mathbf E=i\omega\mu\mathbf H\)、\(\nabla\times\mathbf H=(\sigma-i\omega\varepsilon)\mathbf E\)。消去磁场并代入 \(e^{ikz}\) 平面波，得到 \(k^2=\mu\varepsilon\omega^2+i\mu\sigma\omega\)，虚部代表导电损耗。
(2) 代入 \(k=\beta+i\alpha\) 得 \(\beta^2-\alpha^2=\mu\varepsilon\omega^2\)、\(2\alpha\beta=\mu\sigma\omega\)。取向右传播并衰减的正根，\(\beta=\omega\sqrt{\mu\varepsilon/2}\sqrt{\sqrt{1+(\sigma/(\omega\varepsilon))^2}+1}\)，\(\alpha=\omega\sqrt{\mu\varepsilon/2}\sqrt{\sqrt{1+(\sigma/(\omega\varepsilon))^2}-1}\)。
(3) 在 \(\sigma/(\omega\varepsilon)\gg1\) 时，\(\alpha\simeq\beta\simeq\sqrt{\mu\sigma\omega/2}\)。场振幅 \(e^{-\alpha z}\) 的 \(1/e\) 穿透深度为 \(\delta=1/\alpha=\sqrt{2/(\mu\sigma\omega)}\)，相速度 \(v_p=\omega/\beta\simeq\sqrt{2\omega/(\mu\sigma)}\)。
(4) 时间平均坡印廷流由电场与磁场幅值的乘积组成，因此按 \(e^{-2\alpha z}\) 而非 \(e^{-\alpha z}\) 衰减，功率的 \(1/e\) 长度是 \(\delta/2\)。当 \(\sigma\to0\) 时，\(\alpha\to0\)、\(\beta\to\omega\sqrt{\mu\varepsilon}\)，恢复无损介质波速 \(1/\sqrt{\mu\varepsilon}\)。
